\documentclass{../ElectromagneticFieldTemplate}

\graphicspath{{./Figures/}{../TemplateFigures/}}

% ==================== 封面信息（请根据实际情况填写）====================
\setExperimentInfo
    {电磁波参量的测量}
    {2024217802}
    {严子竣}
    {2024212622}
    {张璇}
    {\today}              % 自动生成当前日期，也可手动写“2026年5月”

% ==================== 正文开始 ====================
\begin{document}
\maketitle

\section{实验目的}
\begin{enumerate}
    \item 在学习均匀平面电磁波特性的基础上，观察电磁波传播特性。
    \item 熟悉并利用相干波原理，测定自由空间内电磁波波长，并确定相位常数和波速。
\end{enumerate}

\section{实验原理}
两束等幅、同频率的均匀平面电磁波在自由空间内以相反方向传播时，会发生干涉现象，形成驻波场分布。本实验利用相干波原理，通过测定驻波场节点的位置变化，求得自由空间内电磁波波长 \(\lambda\)，进而得到相位常数 \(\beta\) 和波速 \(\nu\)。

如图1-1所示（参见实验讲义），入射波经介质板分束后分别被固定金属板 \(P_{r1}\) 和可动金属板 \(P_{r2}\) 反射，两束反射波在接收喇叭 \(P_{r3}\) 处叠加。移动 \(P_{r2}\) 改变两束波的波程差，当波程差满足一定条件时，合成场出现极小值（节点）。相邻节点对应的波程差变化为半个波长，通过测量多个节点的位置，可按下式计算波长平均值：
\[
\lambda = \frac{2(L_N - L_0)}{N}
\]
其中 \(L_0\) 为第一个节点位置，\(L_N\) 为第 \(N+1\) 个节点位置，\(N\) 为半波长个数。进而求得：
\[
\beta = \frac{2\pi}{\lambda}, \qquad \nu = \lambda f
\]
本实验所用信号源中心频率 \(f_0 = 10.2\,\text{GHz}\)，理论波长 \(\lambda_0 = c/f_0 = 29.41\,\text{mm}\)。

\section{实验步骤}
按照讲义中“方法一（直接读零点）”进行测量：
\begin{enumerate}
    \item 整体机械位置初始化：调整发射喇叭 \(P_{r1}\) 与接收喇叭 \(P_{r3}\) 使长边平行于地面，极化一致，且接收喇叭轴向与固定板 \(P_{r1}\) 轴向重合。
    \item 按图1-1安装铝板 \(P_{r1}\)、\(P_{r2}\) 及有机玻璃板，移动板 \(P_{r2}\) 安装在螺旋测微器上。使电磁波以 \(45^\circ\) 角入射到有机玻璃板，再垂直入射到 \(P_{r1}\) 与 \(P_{r2}\)。
    \item 锁紧各部件，调整发射功率使接收喇叭示波器上信号峰峰值最大值约为1V。
    \item 旋转螺旋测微器，将可动反射板 \(P_{r2}\) 移至最右侧（起始位置），然后缓慢向左移动，观察示波器上信号峰峰值变化。
    \item 当信号峰峰值出现极小值时，从螺旋测微器上记录该节点位置 \(L_0\)。继续移动 \(P_{r2}\)，依次记录后续 \(N\) 个极小值点位置 \(L_1, L_2, \ldots, L_N\)。本次实验取 \(N=3\)，共记录4个节点位置。
    \item 将测得数据代入公式计算波长 \(\lambda\)、相位常数 \(\beta\) 和波速 \(\nu\)。
\end{enumerate}

\section{实验数据}
螺旋测微器读数（节点位置）记录如下：
\begin{table}[H]
    \centering
    \caption{节点位置测量数据}
    \begin{tabular}{c|c}
        \hline
        节点编号 & 位置 \(L\) (mm) \\
        \hline
        \(L_0\) & 65.75 \\
        \(L_1\) & 47.06 \\
        \(L_2\) & 34.77 \\
        \(L_3\) & 19.82 \\
        \hline
    \end{tabular}
\end{table}
信号源频率 \(f_0 = 10.2\,\text{GHz}\)，理论波长 \(\lambda_0 = 29.41\,\text{mm}\)。

\section{实验结果整理与误差分析}
\subsection{计算结果}
取 \(N=3\)（4个节点），代入公式：
\[
\lambda = \frac{2(L_3 - L_0)}{N} = \frac{2 \times |19.82 - 65.75|}{3}\,\text{mm} = \frac{2 \times 45.93}{3}\,\text{mm} = 30.62\,\text{mm}
\]
\[
\beta = \frac{2\pi}{\lambda} = \frac{2\pi}{30.62}\,\text{rad/mm} \approx 0.205\,\text{rad/mm}
\]
\[
\nu = \lambda f = 30.62\times10^{-3}\,\text{m} \times 10.2\times10^9\,\text{Hz} = 3.123\times10^8\,\text{m/s}
\]
文献中自由空间波速 \(c = 2.998\times10^8\,\text{m/s}\)，相对误差：
\[
\frac{|\nu - c|}{c} \times 100\% = \frac{|3.123-2.998|}{2.998}\times100\% \approx 4.17\%
\]

\subsection{误差分析}
\begin{itemize}
    \item \textbf{系统误差}：实验装置中电磁波在近场区传播，并非理想平面波，导致节点间距不完全相等，影响平均波长的精度。
    \item \textbf{读数误差}：螺旋测微器最小分度值为0.01 mm，但人眼判断“极小值点”存在主观性，可能引入读数偏差。
    \item \textbf{环境干扰}：周围金属物体对电磁波的反射可能造成场分布畸变，使节点位置偏移。
    \item \textbf{极化与对准误差}：收发喇叭的极化方向、轴线对准若不严格，也会影响驻波比和节点位置的准确判定。
\end{itemize}

\subsection{讨论}
所测波长 \(\lambda = 30.62\,\text{mm}\) 较理论值 \(\lambda_0 = 29.41\,\text{mm}\) 偏大约4\%，主要原因是测量时节点位置受近场效应影响，尤其是靠近发射喇叭的节点位置偏差较大。通过取多个节点平均（\(N=3\)）已在很大程度上减小了随机误差。若增加测量节点数（如 \(N=4\) 或更大），可进一步提高测量精度。

\section{思考题}
\subsection{相干法测电磁波波长时，图1-1中介质板放置位置若转 \(90^\circ\)，将出现什么现象？这时能否测准 \(\lambda\)？为什么？}
发射端发射到介质板处的反射波直接进入接收端\(P_{r3}\)，透射波经\(P_{r2}\)反射，反射波、折射波均未被接收端接收到，不会形成驻波，无法进行测量。

\subsection{通过本实验，你对金属波导有哪些了解？}
本实验中的发射和接收喇叭属于波导结构的一部分（矩形波导馈电）。金属波导能够引导电磁波定向传播，具有低损耗、高功率容量等优点。波导的尺寸决定了其截止频率，只有当工作频率高于截止频率时，波导才能有效传输电磁波。

\end{document}